% Chapter 1: Introduction to Embedded Rust and Project Architecture
% File: chapters/chapter-01-introduction.tex

\chapter{Introduction to Embedded Rust and Project Architecture}
\label{ch:chapter-01}
\chapthumb{1}

\begin{quote}
    \textit{"The biggest memory safety vulnerability is the assumption that memory safety isn't important for embedded systems."}
    \\
    -- Embedded Systems Security Research, 2023
\end{quote}

This chapter introduces embedded Rust development and establishes the foundation for understanding the PAC/HAL/BSP architecture pattern. We'll explore why Rust is compelling for embedded systems, create a bare-metal program, and set up the complete development environment.

\minitoc

\clearpage

% ============================================================================
% SECTION 01.01: Why Rust for Embedded Systems
% ============================================================================
% Section 1.1: Why Rust for Embedded Systems

Complex software systems, especially those dealing with direct hardware manipulation, memory management, and concurrent interrupts, are notoriously difficult to write correctly. Traditional embedded development often relies on C or C++, languages that provide raw power and direct hardware access but lack modern safety guarantees. Memory safety bugs---buffer overflows, use-after-free, null pointer dereferences, and data races---account for the majority of security vulnerabilities and system crashes in embedded systems.

\section{Rust's Value Proposition for Embedded}
\label{sec:why-rust-value-prop}

Rust addresses these challenges through a unique combination of features that make it exceptionally well-suited for embedded systems programming:

\subsection{Memory Safety Without Garbage Collection}

Traditional memory-safe languages rely on garbage collection, which introduces unpredictable pauses and significant runtime overhead---unacceptable constraints for resource-constrained microcontrollers. Rust achieves memory safety through its ownership system, enforced at compile time with zero runtime cost.

\begin{notebox}
    \textbf{Key Insight}: Rust's borrow checker prevents memory safety bugs at compilation, eliminating entire classes of vulnerabilities before the code ever reaches the hardware.
\end{notebox}

Consider this common C bug pattern:

\begin{minted}{c}
void process_data(uint8_t *buffer, size_t len) {
    uint8_t local_buffer[16];
    // Buffer overflow if len > 16!
    memcpy(local_buffer, buffer, len);
    // ... process data
}
\end{minted}

The same logic in Rust:

\begin{minted}{rust}
fn process_data(buffer: &[u8]) {
    // Compile error: attempt to copy 32 bytes into 16-byte array
    let mut local_buffer = [0u8; 16];
    local_buffer.copy_from_slice(buffer); // Panics if buffer.len() != 16
    // Or use safe slice copying:
    let count = buffer.len().min(16);
    local_buffer[..count].copy_from_slice(&buffer[..count]);
}
\end{minted}

The Rust compiler enforces bounds checking at compile time or generates explicit panics for runtime bounds violations, preventing undefined behavior.

\subsection{Zero-Cost Abstractions}

Rust's high-level abstractions compile down to the same efficient machine code as hand-written C. This is critical for embedded systems where every byte and CPU cycle matters.

\begin{example}
A GPIO pin abstraction in Rust:

\begin{minted}{rust}
let pin = gpio.pin0.into_output();
pin.set_high();
\end{minted}

This compiles to a single register write instruction, identical to:

\begin{minted}{c}
GPIO->DOUT |= (1 << 0);
\end{minted}
\end{example}

The Rust code provides type safety, compile-time mode checking, and clear semantics while maintaining identical performance to the raw C version.

\subsection{Fearless Concurrency}

Embedded systems increasingly require concurrent operation: interrupt handlers, DMA transfers, and multiple sensors polling simultaneously. Rust's ownership model prevents data races at compile time.

\begin{minted}{rust}
// This compiles safely
static mut COUNTER: u32 = 0;

#[interrupt]
fn TIMER0() {
    // Safe because we use atomic operations
    cortex_m::interrupt::free(|_| {
        unsafe { COUNTER += 1; }
    });
}

// This fails to compile!
fn increment_counter() {
    COUNTER += 1; // Error: cannot access mutable static without unsafe
}
\end{minted}

The compiler forces developers to explicitly handle synchronization, making race conditions nearly impossible.

\subsection{Comparison with Traditional Embedded Languages}
\label{sec:language-comparison}

\subsection{Rust vs C}

\begin{table}[H]
    \centering
    \begin{tabular}{l c c}
        \toprule
        \textbf{Feature} & \textbf{C} & \textbf{Rust} \\
        \midrule
        Memory Safety & Manual (runtime errors) & Compile-time guaranteed \\
        Concurrency Safety & Manual (locks, careful design) & Compile-time guaranteed \\
        Type Safety & Weak (void*, casts) & Strong (no implicit casts) \\
        Null Pointers & Always possible & Option<T> type \\
        Error Handling & Error codes & Result<T, E> type \\
        Abstraction Cost & Potentially high & Zero-cost \\
        Learning Curve & Moderate & Steep (ownership concepts) \\
        Ecosystem & Mature, established & Growing rapidly \\
        Real-time Suitability & Excellent (predictable) & Excellent (no GC) \\
        \bottomrule
    \end{tabular}
    \caption{Rust vs C for embedded development}
    \label{tab:rust-vs-c}
\end{table}

\subsection{Real-World Impact}

The Linux kernel, which has historically been written entirely in C, began accepting Rust code in version 6.1 (2022). This demonstrates Rust's maturity for systems programming even at the highest levels.

In the embedded space, companies are reporting:

\begin{itemize}
    \item 40-50% reduction in debugging time
    \item Near elimination of memory corruption bugs
    \item Improved team productivity through better compiler diagnostics
    \item Easier onboarding for new developers
\end{itemize}

\subsection{Rust's Embedded Ecosystem}
\label{sec:rust-embedded-ecosystem}

\subsection{Core Crates}

The embedded Rust ecosystem provides standardized interfaces:

\begin{itemize}
    \item \crate{embedded-hal}: Hardware abstraction layer traits for portability
    \item \crate{cortex-m}: Low-level Cortex-M register access
    \item \crate{cortex-m-rt}: Startup code and interrupt handling
    \item \crate{cortex-m-semihosting}: Debug output via debugger
    \item \crate{svd2rust}: PAC generation from SVD files
    \item \crate{panic-halt}: Panic handler for minimal footprint
\end{itemize}

\subsection{Tooling Integration}

Rust integrates seamlessly with embedded development workflows:

\begin{itemize}
    \item \textbf{cargo}: Built-in package manager and build system
    \item \textbf{probe-rs}: Flash programming and debugging
    \item \textbf{defmt}: Efficient logging for resource-constrained targets
    \item \textbf{flip-link}: Linker wrapper for stack overflow protection
\end{itemize}

\subsection{Addressing Embedded Development Challenges}
\label{sec:embedded-challenges}

\subsection{Resource Constraints}

Rust's ownership system enables compile-time memory management without runtime overhead. The \codeinline{no_std} environment runs without heap allocation, relying entirely on stack allocation and static data.

\begin{minted}{rust}
#![no_std]  // No standard library (no heap allocation)
#![no_main] // Disabled main wrapping

// All memory is statically allocated
static BUFFER: [u8; 1024] = [0; 1024];

#[no_mangle]
pub extern "C" fn main() -> ! {
    loop {
        // Application logic
    }
}
\end{minted}

\subsection{Deterministic Behavior}

Embedded systems require predictable timing. Rust provides:

\begin{itemize}
    \item No hidden allocations or deallocations
    \item No garbage collection pauses
    \item Explicit control over memory layout
    \item \codeinline{const fn} for compile-time computation
\end{itemize}

\subsection{Hardware Interface Safety}

Direct hardware register manipulation is inherently unsafe. Rust forces developers to mark unsafe operations explicitly:

\begin{minted}{rust}
pub fn enable_gpio_clock() {
    unsafe {
        (*CMU::ptr()).hfbusclken0.modify(|_, w|
            w.gpio().set_bit()
        );
    }
}
\end{minted}

The \codeinline{unsafe} block signals to developers and reviewers that manual verification is required.

\subsection{The EFR32MG24: Ideal for Rust}
\label{sec:efr32mg24-for-rust}

The Silicon Labs EFR32MG24 is particularly well-suited for Rust development:

\begin{itemize}
    \item \textbf{ARM Cortex-M33}: Modern architecture with TrustZone-M
    \item \textbf{Memory Protection Unit}: Hardware-enforced memory isolation
    \item \textbf{Deterministic Architecture}: No unpredictable timing behavior
    \item \textbf{Rich Peripherals}: 133 peripherals across multiple domains
    \item \textbf{Unified Radio}: Multi-protocol stack support
    \item \textbf{Hardware Cryptography}: Secure element integration
\end{itemize}

\subsection{Summary}
\label{sec:why-rust-summary}

Rust brings modern software engineering practices to embedded development without sacrificing performance or control. Its compile-time safety guarantees eliminate entire classes of bugs that plague C-based embedded systems, while its zero-cost abstractions enable expressive, maintainable code that compiles to efficient machine instructions.

The journey through this book will demonstrate:

\begin{enumerate}
    \item How to generate a Peripheral Access Crate (PAC) from vendor SVD files
    \item Building safe, ergonomic HAL abstractions on top of raw registers
    \item Implementing standard embedded-hal traits for ecosystem compatibility
    \item Creating board support packages for specific hardware
    \item Integrating complex protocol stacks (Bluetooth, Zigbee, Thread)
    \item Testing strategies for embedded systems
    \item Production deployment and maintenance
\end{enumerate}

Along the way, we'll explore the unique challenges and solutions that emerged during the real-world implementation of EFR32MG24 support in Rust.
\clearpage


% ============================================================================
% SECTION 01.02: Your First Bare-Metal Rust Program
% ============================================================================
\section{Your First Bare-Metal Rust Program}
\label{sec:bare-metal-program}

% Section 01.02: Your First Bare-Metal Rust Program

After understanding why Rust is compelling for embedded systems, let's create our first bare-metal Rust program. This example will run without an operating system, directly on the microcontroller hardware.

\subsection{Creating a Minimal Embedded Project}

First, create a new directory for our project:

\begin{minted}{bash}
$ mkdir blink_bare_metal
$ cd blink_bare_metal
$ cargo init --name blink_bare_metal
\end{minted}

This creates a standard Rust project structure. However, for embedded development, we need to modify it significantly.

\subsection{Configuring for bare\_metal: \texttt{no\_std}}

Before we begin, let's understand what a bare-metal embedded project looks like. Edit \filepath{Cargo.toml}:

\begin{minted}{toml}
[package]
name = "blink_bare_metal"
version = "0.1.0"
edition = "2021"

[profile.dev]
opt-level = "z"          # Optimize for size
debug = true

[profile.release]
opt-level = "z"          # Optimize for size
panic = "abort"          # Simplest panic handler

[dependencies]
# We'll add dependencies as needed
\end{minted}

Now modify \filepath{src/main.rs}:

\begin{minted}{rust}
#![no_std]    // Don't link the standard library (no heap, no OS)
#![no_main]   // We define our own entry point

// Prevent Rust from generating main() wrapper
use core::panic::PanicInfo;

/// Our entry point - Cortex-M calls this on reset
#[no_mangle]
pub extern "C" fn main() -> ! {
    // Application logic will go here

    loop {
        // Infinite loop to prevent return
    }
}

/// Simple panic handler - just halt
#[panic_handler]
fn panic(_info: &PanicInfo) -> ! {
    loop {}
}
\end{minted}

\subsection{What is no\_std?}

\begin{notebox}
    \textbf{no\_std} means ``no standard library.'' Our program cannot use:
    \begin{itemize}
        \item Heap allocation (Box, Vec, String)
        \item File I/O
        \item Threads
        \item Networking
    \end{itemize}

    We can still use: \codeinline{core}, slices, static arrays, and direct hardware access.
\end{notebox}

The standard library requires an operating system for features like heap allocation and threads. In bare-metal embedded, we have no OS, so we use only \codeinline{core}, which provides fundamental types and operations that don't require an OS.

\subsection{The Entry Point}

In hosted environments (with an OS), Rust generates a C main() function that calls your main(). In bare metal:

\begin{itemize}
    \item \codeinline{#[no_mangle]} prevents Rust from changing the function name
    \item \codeinline{pub extern "C" fn} ensures the correct calling convention
    \item \codeinline{-> !} means "never returns" (diverging function)
    \item We must provide an infinite loop or halt instruction
\end{itemize}

\subsection{Cross-Compilation}

Our development machine (likely x86\_64) can't execute ARM Cortex-M code. We need to cross-compile:

\begin{minted}{bash}
$ rustup target add thumbv8m.main-none-eabihf
\end{minted}

This installs the Rust standard library for the ARM Cortex-M33 architecture without an OS.

Create \filepath{.cargo/config.toml}:

\begin{minted}{toml}
[build]
target = "thumbv8m.main-none-eabihf"

[target.thumbv8m.main-none-eabihf]
# No `runner` on purpose: `cargo run` would program the chip the moment it builds,
# and on this board that has to go through the guarded script (see the note below).
rustflags = [
  "-C", "link-arg=-Tlink.x",
  "-C", "link-arg=--nmagic",
]

[env]
DEFMT_LOG = "debug"
\end{minted}

\begin{quote}
\textbf{Why this book has no cargo run.} The first 24 KB of flash
belong to the Gecko bootloader. A `runner` in `.cargo/config.toml` turns every
\codeinline{cargo run} into a write, and a single wrong linker script silently relocates
the vector table into that region. This project flashes through
\codeinline{scripts/flash-lnx.sh}: dry run by default, \codeinline{--confirm} to write,
a hard refusal when the image's lowest address is below \codeinline{0x08006000}, and a
read-back verification afterwards. The same script is what re-flashes the board after a
mistake — which matters more than the convenience of \codeinline{cargo run}.
\end{quote}

\subsection{Building the Project}

\begin{minted}{bash}
$ cargo build --release
   Compiling blink_bare_metal v0.1.0
    Finished release [optimized] target(s) in 0.5s
\end{minted}

Check the result:

\begin{minted}{bash}
$ cargo size --release
   text    data     bss     dec     hex filename
     64       0       0      64      40 blink_bare_metal
\end{minted}

Our minimal program is only 64 bytes!

\subsection{Analyzing the Binary}

Use \codeinline{objdump} to examine the machine code:

\begin{minted}{bash}
$ cargo objdump --release -- -d
\end{minted}

This shows the ARM assembly instructions. You'll see something like:

\begin{minted}{text}
bare_metal:     file format elf32-littlearm

Disassembly of section .text:

main:
    b #-4 <main>
    b #-4 <main>
\end{minted}

The two \codeinline{b #-4} instructions are branches to themselves---infinite loops. One is our main() function, the other is the panic handler.

\subsection{What This Demonstrates}

Even at this minimal stage, we've established:

\begin{itemize}
    \item Cross-compilation toolchain setup
    \item Bare-metal runtime configuration
    \item Minimal program that builds successfully
    \item Linker script and startup code (provided by cortex-m-rt)
\end{itemize}

The program won't do anything visible yet (no LEDs blinking), but it represents a working foundation. In the next section, we'll add actual hardware interaction.

\begin{exercise}[Easy]
    Modify the panic handler to increment a counter in an infinite loop. Use a static mut variable and examine how the compiler warns about unsafe code.
\end{exercise}

\begin{exercise}[Medium]
    Research what happens when main() (or the reset handler) returns on a Cortex-M microcontroller. Why is \codeinline{-> !} important?
\end{exercise}

\begin{exercise}[Hard]
    Create a minimal linker script (\filepath{linker.ld}) and build script (\filepath{build.rs}) to understand how the linker allocates memory for your program.
\end{exercise}
\clearpage


% ============================================================================
% SECTION 01.03: The PAC/HAL/BSP Architecture Pattern
% ============================================================================
\section{The PAC/HAL/BSP Architecture Pattern}
\label{sec:pac-hal-bsp-architecture}

% Section 01.03: The PAC/HAL/BSP Architecture Pattern

Embedded Rust development follows a layered architecture that promotes code reuse, safety, and separation of concerns. This pattern is not arbitrary—it reflects years of lessons learned from implementing hardware support across multiple vendors and microcontroller families.

\subsection{The Three-Layer Architecture}
\label{subsec:three-layer-architecture}

\begin{figure}[H]
    \centering
    \begin{tikzpicture}[
        layer/.style={rectangle, draw=black, thick, minimum width=8cm, minimum height=1.5cm, font=\bfseries},
        arrow/.style={-Latex, thick, double},
        label/.style={font=\small\itshape}
    ]

        % Application Layer
        \node[layer, fill=rustblue!20] (app) at (0, 4.5) {Application Code};

        % BSP Layer
        \node[layer, fill=green!20] (bsp) at (0, 3) {Board Support Package};

        % HAL Layer
        \node[layer, fill=yellow!20] (hal) at (0, 1.5) {Hardware Abstraction Layer};

        % PAC Layer
        \node[layer, fill=red!20] (pac) at (0, 0) {Peripheral Access Crate};

        % Hardware Layer
        \node[layer, fill=gray!20] (hw) at (0, -1.5) {Hardware Peripherals};

        % Arrows showing dependency
        \draw[arrow] (app) -- (bsp);
        \draw[arrow] (bsp) -- (hal);
        \draw[arrow] (hal) -- (pac);
        \draw[arrow] (pac) -- (hw);

        % Labels
        \node[label, left=2cm of app] {Board-Specific};
        \node[label, left=2cm of bsp] {Portable};
        \node[label, left=2cm of hal] {Register Access};
        \node[label, left=2cm of pac] {Hardware};

        % Portability indicators
        \node[font=\small, right=2cm of app] {Application-Specific};
        \node[font=\small, right=2cm of bsp] {Board-Specific};
        \node[font=\small, right=2cm of hal] {Chip Family Specific};
        \node[font=\small, right=2cm of pac] {Vendor Provided / Generated};
        \node[font=\small, right=2cm of hw] {Fixed by Hardware};
    \end{tikzpicture}
    \caption{The PAC/HAL/BSP layered architecture}
    \label{fig:pac-hal-bsp-layers}
\end{figure}

Each layer has distinct responsibilities and portability characteristics.

\subsection{Peripheral Access Crate (PAC)}
\label{subsec:pac-layer}

The PAC provides the lowest-level, register-accurate representation of the microcontroller's peripherals. It is typically generated automatically from the vendor's System View Description (SVD) file using \codeinline{svd2rust}.

\textbf{Key characteristics}:
\begin{itemize}
    \item One-to-one mapping with hardware registers
    \item No abstractions on top of hardware
    \item All operations are \codeinline{unsafe} (by design)
    \item Generated automatically—never hand-edited
    \item Chip variant specific (EFR32MG24B220 vs A020)
\end{itemize}

\begin{minted}{rust}
// Example: Direct register access via PAC
unsafe {
    (*efr32mg24_pac::GPIO_S::ptr())
        .port_a_dout
        .modify(|r, w| w.bits(r.bits() | (1 << 5)));
}
\end{minted}

This writes directly to GPIO Port A, bit 5. The \codeinline{unsafe} keyword forces you to acknowledge that you're bypassing Rust's safety guarantees.

\begin{warningbox}
    PAC code should be wrapped in HAL abstractions. Direct PAC usage in application code is an anti-pattern that sacrifices safety for no benefit.
\end{warningbox}

\subsection{Hardware Abstraction Layer (HAL)}
\label{subsec:hal-layer}

The HAL provides safe, ergonomic APIs that wrap the raw PAC access. It implements standard traits for portability and adds compile-time safety guarantees.

\textbf{Key characteristics}:
\begin{itemize}
    \item Safe APIs (no \codeinline{unsafe} required by users)
    \item Type-safe state machines (e.g., pin modes encoded in types)
    \item Implements \crate{embedded-hal} traits
    \item Chip family specific (works across all EFR32MG24 variants)
    \item Medium abstraction level
    \item Written and maintained by community
\end{itemize}

\begin{minted}{rust}
// Example: GPIO via HAL
let pin = gpio.pin5.into_output();
pin.set_high();
\end{minted}

Compare this to the PAC version above. The HAL provides:
\begin{itemize}
    \item Type-safe pin selection (can't use non-existent pins)
    \item Compile-time mode tracking (can't read from output pin)
    \item Clear, readable method names
    \item No unsafe required
    \item Identical performance (zero-cost abstraction)
\end{itemize}

\subsection{Board Support Package (BSP)}
\label{subsec:bsp-layer}

The BSP provides board-specific abstractions and convenience APIs. It's the only layer that "knows" about your specific PCB design.

\textbf{Key characteristics}:
\begin{itemize}
    \item Maps logical functions to physical pins (e.g., LED1 → PA5)
    \item Provides board-specific initialization sequences
    \item Combines peripherals for board features
    \item Board specific (each PCB design has its own BSP)
    \item Highest abstraction level
    \item Minimal code (thin wrappers over HAL)
\end{itemize}

\begin{minted}{rust}
// Example: BSP for XIAO MG24 Sense
use xiao_mg24_bsp::{Board, LED};

let board = Board::new();
board.led_green.on();
delay_ms(500);
board.led_green.off();
\end{minted}

The BSP knows that the green LED connects to PA5 on this particular board and provides a convenient LED abstraction.

\subsection{The Flow of Abstractions}
\label{sec:abstraction-flow}

Here's how a simple operation (turning on an LED) flows through the layers:

\begin{enumerate}
    \item \textbf{Application}: Calls \codeinline{board.led_red.on()}
    \item \textbf{BSP}: Translates to \codeinline{hal.gpio.pin5.set_high()}
    \item \textbf{HAL}: Checks pin mode, then calls PAC
    \item \textbf{PAC}: Executes \codeinline{GPIO->DOUT |= (1 << 5)}
    \item \textbf{Hardware}: Physical pin voltage goes high, LED illuminates
\end{enumerate}

This layered approach provides benefits:

\textbf{Portability}: Switch to a different EFR32MG24 board? Only rewrite the BSP. HAL and PAC remain unchanged.

\textbf{Maintainability}: Bug in GPIO handling? Fix it in the HAL. All boards using that HAL benefit.

\textbf{Composability}: Libraries written for HAL traits work across all boards using that HAL.

\begin{notebox}
    \textbf{Architecture Rule}: Application code should never directly use PAC. Always go through HAL and BSP.
\end{notebox}

\subsection{Why This Separation Matters}
\label{sec:separation-importance}

The embedded Rust community settled on this architecture through experience:

\subsection{PAC Regeneration}

When vendors release updated SVD files (fixing errors or documenting new features), you regenerate the PAC without affecting HAL or BSP code.

\subsection{Multiple HAL Implementations}

Different use cases may need different HAL trade-offs:
\begin{itemize}
    \item High-performance HAL for real-time applications
    \item Small-footprint HAL for resource-constrained devices
    \item Feature-rich HAL for development boards
\end{itemize}

All can use the same PAC.

\subsection{Ecosystem Compatibility}

The \crate{embedded-hal} trait system enables an ecosystem of drivers:
\begin{itemize}
    \item Display drivers (SSD1306, ST7789) work with any SPI implementation
    \item Sensor drivers (BMP280, MPU6050) work with any I2C implementation
    \item No code changes needed when switching microcontrollers
\end{itemize}

\subsection{Our Approach in This Book}
\label{sec:our-approach}

In this book, we'll implement the stack bottom-up:

\begin{enumerate}
    \item \textbf{Chapter 2}: Generate the PAC from Silicon Labs SVD
    \item \textbf{Chapters 4-9}: Build HAL for all peripherals
    \item \textbf{Chapter 11}: RAIL radio abstraction
    \item \textbf{Chapters 12-14}: Protocol stacks (Bluetooth, Zigbee, Thread)
    \item \textbf{Chapter 15}: BSP for XIAO MG24 Sense
\end{enumerate}

This order ensures:
\begin{itemize}
    \item Each layer has a strong foundation
    \item Clear dependency management
    \item Comprehensive testing at each level
    \item Real hardware validation
\end{itemize}

\begin{exercise}[Easy]
    Look at a simple development board (like Arduino Uno) and list:
    \begin{enumerate}
        \item Which features belong in PAC vs HAL vs BSP?
        \item Why does Arduino's "digitalWrite()" combine all three layers?
    \end{enumerate}
\end{exercise}

\begin{exercise}[Medium]
    Design a type-safe state machine for a USART that:
    \begin{enumerate}
        \item Only allows baud rate configuration when USART is disabled
        \item Only allows transmission when USART is enabled
        \item Prevents mode changes during active transmission
    \end{enumerate}
    How would you encode these states in Rust's type system?
\end{exercise}

\begin{exercise}[Hard]
    Research other embedded HAL implementations:
    \begin{enumerate}
        \item Compare \codeinline{stm32-hal} vs \codeinline{nrf-hal}
        \item Note differences in how they handle:
            \begin{itemize}
                \item Clock configuration
                \item Error types
                \item Feature flags
                \item Documentation approach
            \end{itemize}
    \end{enumerate}
\end{exercise}
\clearpage


% ============================================================================
% SECTION 01.04: Project Organization and Repository Structure
% ============================================================================
\section{Project Organization and Repository Structure}
\label{sec:project-organization}

% Section 01.04: Project Organization and Repository Structure

A well-organized repository structure is essential for maintainability, especially in embedded projects where you may support multiple chip variants, boards, and protocol stacks. In this section, we'll examine the repository structure used throughout this book and understand the rationale behind each decision.

\subsection{Top-Level Directory Structure}

Here's the complete structure of our EFR32MG24 project:

% Directory Tree for Section 1.4 - Repository Structure
\begin{figure}[H]
    \centering
    \begin{forest}
        for tree={
            font=\ttfamily\footnotesize,
            grow'=0,
            child anchor=west,
            parent anchor=south,
            anchor=west,
            calign=first,
            edge path={
                \noexpand\path [draw] (!u.south west) +(7.5pt,0) |- node[fill,inner sep=1.25pt] {} (.child anchor);
            },
            before typesetting nodes={
                if n=1
                    {insert before={[,phantom]}}
                    {}
            },
            fit=band,
            before computing xy={l=15pt},
        }
        [EFR32MG24/
            [efr32mg24-pac/
                [Cargo.toml]
                [src/]
                [svd/]
            ]
            [efr32mg24-hal/
                [Cargo.toml]
                [src/]
                [examples/]
            ]
            [efr32mg24-rail-sys/]
            [efr32mg24-bluetooth-sys/]
            [xiao-mg24-bsp/]
            [docs/]
            [book/
                [latex/]
                [TOC.yml]
            ]
            [examples/]
            [Cargo.toml]
            [README.md]
            [CLAUDE.md]
        ]
    \end{forest}
    \caption{EFR32MG24 project repository structure}
    \label{fig:repo-structure}
\end{figure}


\subsection{Crate Organization}

\subsection{Peripheral Access Crate (efr32mg24-pac)}

\begin{minted}{text}
efr32mg24-pac/
├── Cargo.toml          # PAC crate metadata
├── build.rs            # Build script (if needed)
├── svd/
│   └── EFR32MG24B220F1536IM48.svd  # SVD file
├── src/
│   └── lib.rs         # 138,000+ lines of generated code
└── README.md          # PAC-specific documentation
\end{minted}

**Key Points**:
\begin{itemize}
    \item Contains ONLY auto-generated code from svd2rust
    \item Never edited manually after generation
    \item Version locked to specific SVD file version
    \item Re-exported by HAL, users never depend on PAC directly
\end{itemize}

\subsection{Hardware Abstraction Layer (efr32mg24-hal)}

The HAL workspace contains multiple peripheral crates:

\begin{minted}{text}
efr32mg24-hal/
├── Cargo.toml              # Workspace root
├── src/
│   ├── lib.rs             # Main HAL entry point
│   ├── clock/             # Clock management (CMU)
│   │   ├── mod.rs
│   │   ├── types.rs
│   │   ├── traits.rs
│   │   └── README.md
│   ├── gpio/              # GPIO peripheral
│   │   ├── mod.rs
│   │   ├── types.rs
│   │   ├── traits.rs
│   │   └── README.md
│   ├── usart/             # USART communication
│   │   ├── mod.rs
│   │   ├── types.rs
│   │   ├── traits.rs
│   │   └── README.md
│   ├── i2c/               # I2C peripheral
│   │   ├── mod.rs
│   │   ├── types.rs
│   │   ├── traits.rs
│   │   └── README.md
│   ├── spi/               # SPI peripheral
│   │   ├── mod.rs
│   │   ├── types.rs
│   │   ├── traits.rs
│   │   └── README.md
│   └── prelude.rs         # Common imports
├── examples/              # HAL-level examples
│   ├── 01_clock.rs
│   ├── 02_gpio_blink.rs
│   ├── 03_usart_echo.rs
│   └── ...
└── docs/                  # HAL-specific documentation
    ├── README.md
    └── BUILDING.md
\end{minted}

\textbf{File Pattern}: Each peripheral follows the same structure with four files: mod.rs, types.rs, traits.rs, and README.md.

\begin{minted}{text}
peripheral-name/
├── mod.rs       // Main implementation
├── types.rs     // Type definitions
├── traits.rs    // Trait implementations
└── README.md    // Per-peripheral documentation
\end{minted}

\subsection{Radio System Crates (efr32mg24-rail-sys)}

FFI bindings to Silicon Labs RAIL (Radio Abstraction Interface Layer):

\begin{minted}{text}
efr32mg24-rail-sys/
├── Cargo.toml          // Link to vendor libraries
├── build.rs
├── wrapper.h          // bindgen header
└── src/
    └── lib.rs         # Generated FFI bindings
\end{minted}

**Purpose**: Enables multi-protocol radio operation (Bluetooth, Zigbee, Thread, proprietary).

\subsection{Protocol System Crates (bluetooth-sys, etc.)}

FFI bindings to higher-level protocol stacks:

\begin{minted}{text}
efr32mg24-bluetooth-sys/
├── Cargo.toml         # Bluetooth library linking
├── build.rs
├── wrapper.h
└── src/
    ├── lib.rs        # FFI bindings
    └── events/       # Event handling code
\end{minted}

\subsection{Board Support Package (xiao-mg24-bsp)}

Board-specific abstractions:

\begin{minted}{text}
xiao-mg24-bsp/
├── Cargo.toml          # Depends on efr32mg24-hal
├── src/
│   ├── lib.rs         # Board initialization
│   ├── board.rs       # Board struct
│   ├── sensors.rs     # Onboard sensors
│   └── radio.rs       # Radio utilities
└── examples/
    ├── blink_onboard_led.rs
    ├── ble_beacon.rs
    └── zigbee_coordinator.rs
\end{minted}

\subsection{Why This Structure Works}

\subsection{Clear Separation of Concerns}
\begin{itemize}
    \item PAC: Hardware truth (auto-generated)
    \item HAL: Safe abstractions (manually written)
    \item BSP: Board integration (thin wrappers)
    \item Application: User logic (no hardware details)
\end{itemize}

\subsection{Testing Strategy}
Each layer can be tested independently:
\begin{itemize}
    \item PAC: Verify register addresses match datasheet
    \item HAL: Unit tests with mocked registers
    \item BSP: Integration tests on real hardware
    \item Application: Standard Rust testing
\end{itemize}

\subsection{Version Management}
Crates can be versioned independently:
\begin{minted}{text}
efr32mg24-pac = "0.1.0"      # Hardware revision specific
efr32mg24-hal = "0.4.2"      # Feature updates
xiao-mg24-bsp = "0.2.1"      # Board revision specific
\end{minted}

\begin{exercise}[Easy]
    Sketch a project structure for supporting three boards:
    \begin{itemize}
        \item Seeed XIAO MG24 (covered here)
        \item Silicon Labs xG24-DK2601B dev kit
        \item Custom PCB with external sensors
    \end{itemize}
    What crates would be shared? Which would be board-specific?
\end{exercise}

\begin{exercise}[Medium]
    Consider a startup shipping products using this architecture:
    \begin{enumerate}
        \item How would you version the crates when releasing new firmware?
        \item What testing strategy ensures HAL changes don't break BSP?
        \item How would you handle multiple hardware revisions?
    \end{enumerate}
\end{exercise}

\begin{exercise}[Hard]
    Compare this structure to Linux kernel architecture:
    \begin{enumerate}
        \item Map PAC/HAL/BSP to Linux equivalents
        \item How does Linux handle different architectures? (arch/ directory)
        \item How does Linux handle different boards? (Device Tree)
        \item What lessons can we apply to our Rust embedded architecture?
    \end{enumerate}
\end{exercise}
\clearpage


% ============================================================================
% SECTION 01.05: Development Prerequisites and Toolchain Setup
% ============================================================================
\section{Development Prerequisites and Toolchain Setup}
\label{sec:development-prerequisites}

% Section 01.05: Development Prerequisites and Toolchain Setup

Before we can build and flash embedded Rust code to the EFR32MG24, we need to install and configure several tools. This section provides a complete setup guide for Linux, macOS, and Windows.

\subsection{Required Tools}
\label{sec:required-tools}

\textbf{Rust Toolchain}:

\begin{itemize}
    \item \textbf{Rust 1.70+} with support for the thumbv8m.main-none-eabihf target
    \item \textbf{cargo-binutils} for examining compiled binaries
    \item \textbf{cargo-size} for checking binary size
    \item \textbf{cargo-objdump} for disassembly
\end{itemize}
\subsection{Installation on Linux (Ubuntu/Debian)}
\label{subsec:install-linux}
\begin{minted}{bash}
# Install Rust via rustup
curl --proto '=https' --tlsv1.2 -sSf https://sh.rustup.rs | sh

source $HOME/.cargo/env

# Add ARM Cortex-M33 target
rustup target add thumbv8m.main-none-eabihf

# Install cargo tools
cargo install cargo-binutils
rustup component add llvm-tools-preview

# Install system dependencies
sudo apt-get update
sudo apt-get install -y \
    build-essential \
    pkg-config \
    libssl-dev \
    git \
    wget
\end{minted}
\subsection{Installation on macOS}
\label{subsec:install-macos}
\begin{minted}{bash}
# Install Rust via rustup
curl --proto '=https' --tlsv1.2 -sSf https://sh.rustup.rs | sh

source $HOME/.cargo/env

# Add ARM Cortex-M33 target
rustup target add thumbv8m.main-none-eabihf

# Install cargo tools
cargo install cargo-binutils
rustup component add llvm-tools-preview

# Install Homebrew if not already installed
/bin/bash -c "$(curl -fsSL https://raw.githubusercontent.com/Homebrew/install/HEAD/install.sh)"

# Install system dependencies
brew install git wget
\end{minted}
\subsection{Installation on Windows}
\label{subsec:install-windows}
\begin{enumerate}
    \item Download and install from https://rustup.rs/
    \item Run the installer with default settings
    \item After installation, open a new PowerShell/Command Prompt
    \item Install the ARM target:
\end{enumerate}

\begin{minted}{bash}
rustup target add thumbv8m.main-none-eabihf
cargo install cargo-binutils
rustup component add llvm-tools-preview
\end{minted}

\subsection{Hardware Debug Probe Setup}
\label{sec:debug-probe-setup}

The EFR32MG24 requires a debug probe for flashing and debugging. We'll use probe-rs, a Rust-based tool.

\subsection{Installing probe-rs}
\label{subsec:install-probe-rs}
\begin{minted}{bash}
# Install probe-rs
cargo install probe-rs --features cli

# Verify installation
probe-rs list
\end{minted}

Expected output (when hardware is connected):

\begin{minted}{text}
Probing OpenUSB... Done
[0]: J-Link Silicon Labs (<serial-number>) (J-Link)
\end{minted}

\subsection{Flashing: one sanctioned path}
\label{subsec:install-silabs-tools}

Silicon Labs' Simplicity Commander exists, and it is the right tool for one job: the
authorisation that unlocks a secured part. It is \emph{not} the flashing path this book
uses, because a bare \codeinline{commander flash} will happily write wherever the image
says — including over the bootloader.

What this project uses instead, and why it is worth the extra script:

\begin{minted}{bash}
# 1. show where the image would land, without touching the board
python3 scripts/elf-load-range.py target/thumbv8m.main-none-eabihf/release/examples/pac_imu
# load_min=0x08006000 load_max=0x08006813 size=2068 segments=3

# 2. dry run: prints the plan and refuses anything with load_min < 0x08006000
bash scripts/flash-lnx.sh "$ELF"

# 3. write, then verify by reading the flash back
bash scripts/flash-lnx.sh "$ELF" --confirm
\end{minted}

The board hangs off a separate host (\codeinline{lnx}) reachable over SSH, which is why
the script copies the ELF over and drives the programmer there. If your board is on the
machine you build on, the same script works locally — the guards, not the transport, are
the point.

\begin{quote}
\textbf{Measured caveat.} \codeinline{probe-rs} works for the first attach and then
desynchronises on this board's onboard CMSIS-DAP bridge
(\codeinline{Command ID in response (0x14) does not match sent command ID}); a USB
reset does not clear it and only a physical replug does. OpenOCD talks to the same probe
without the problem, so the script uses OpenOCD for reads and \codeinline{probe-rs} for
the write-and-verify step.
\end{quote}

\subsection{Updating Debug Probe Firmware}
\label{subsec:update-firmware}

If using a J-Link probe:

\begin{minted}{bash}
# On Linux/macOS
JLinkExe -CommanderScript <(echo "exit")

# Follow prompts to update firmware
\end{minted}

\subsection{Development Environment Setup}
\label{sec:dev-env-setup}

\subsection{VS Code Configuration}
\label{subsec:vscode-setup}

Install these extensions:
\begin{itemize}
    \item rust-analyzer (official Rust support)
    \item Even Better TOML (for Cargo.toml files)
    \item CodeLLDB (for debugging)
\end{itemize}

Create \filepath{.vscode/settings.json}:
\begin{minted}{json}
{
    "rust-analyzer.cargo.target": "thumbv8m.main-none-eabihf",
    "rust-analyzer.checkOnSave.allTargets": false,
    "rust-analyzer.cargo.features": ["rt"],
    "rust-analyzer.procMacro.ignored": {
        "cortex-m-rt": ["interrupt"]
    }
}
\end{minted}

\subsection{Cargo Configuration for Embedded}
\label{subsec:cargo-config}

Create \filepath{.cargo/config.toml}:
\begin{minted}{toml}
[target.thumbv8m.main-none-eabihf]
rustflags = [
    "-C", "link-arg=-Tlink.x",
    "-C", "link-arg=-Tdefmt.x",  # If using defmt
]

runner = "probe-rs run --chip EFR32MG24B220F1536IM48"

[build]
target = "thumbv8m.main-none-eabihf"

[env]
DEFMT_LOG = "debug"
\end{minted}

\subsection{Git Configuration}
\label{subsec:git-config}

Set up Git for embedded development:
\begin{minted}{bash}
git config --global user.name "Your Name"
git config --global user.email "your.email@example.com"

# Set up LFS for large files (like SVD files)
git lfs install
\end{minted}

Create \filepath{.gitattributes}:
\begin{minted}{text}
# Track large files with Git LFS
*.svd filter=lfs diff=lfs merge=lfs -text
*.pack filter=lfs diff=lfs merge=lfs -text
*.a filter=lfs diff=lfs merge=lfs -text
\end{minted}

\subsection{Verifying Your Setup}
\label{sec:verify-setup}

Run this verification script to ensure everything is correctly installed:
\begin{minted}{bash}
#!/bin/bash

echo "=== Checking Rust Installation ==="
rustc --version
cargo --version

echo ""
echo "=== Checking Targets ==="
rustup target list | grep "thumbv8m.main-none-eabihf (installed)"

echo ""
echo "=== Checking Cargo Tools ==="
cargo size --version 2>/dev/null && echo "✓ cargo-size installed"

echo ""
echo "=== Checking probe-rs ==="
probe-rs --version

echo ""
echo "=== Checking debug probe ==="
probe-rs list
\end{minted}

Expected output:
\begin{minted}{text}
=== Checking Rust Installation ===
rustc 1.70.0 (90c541806 2023-05-31)
cargo 1.70.0 (90c541806 2023-05-31)

=== Checking Targets ===
thumbv8m.main-none-eabihf (installed)

=== Checking Cargo Tools ===
✓ cargo-size installed

=== Checking probe-rs ===
probe-rs 0.18.0

=== Checking debug probe ===
[0]: J-Link Silicon Labs (xxxx) (J-Link)
\end{minted}

If all checks pass, you're ready to start development!

\begin{tipbox}
    \textbf{Common Issues}:
    \\• probe-rs list shows nothing: Check USB cable, ensure probe has power
    \\• Target not installed: Run rustup target add thumbv8m.main-none-eabihf
    \\• cargo-size not found: Run cargo install cargo-binutils
    \\• Permission denied (Linux): Add your user to dialout or plugdev group
\end{tipbox}

\begin{exercise}[Easy]
    Create a shell script that automates the entire setup verification process for your team.
\end{exercise}

\begin{exercise}[Medium]
    Write a cargo subcommand that checks for required tools and provides helpful installation instructions if anything is missing.
\end{exercise}

\begin{exercise}[Hard]
    Using Docker, create a container with the complete development environment pre-installed so new team members can start immediately.
\end{exercise}
\clearpage


% ============================================================================
% SECTION 01.06: Building Your First Embedded Binary
% ============================================================================
\section{Building Your First Embedded Binary}
\label{sec:building-first-binary}

% Section 01.06: Building Your First Embedded Binary

Now that we have our toolchain set up and understand the project structure, let's build and analyze a real embedded binary. This section will guide you through the complete build process with detailed explanations of each step.

\subsection{Preparing a Complete Example}
\label{sec:complete-example}

Let's create a minimal but complete embedded project that we can actually build and analyze:
\begin{minted}{bash}
$ cargo new --lib efr32mg24-minimal
$ cd efr32mg24-minimal
$ mkdir examples
$ cat > examples/minimal.rs << 'EOF'
#![no_std]
#![no_main]

use core::panic::PanicInfo;

#[panic_handler]
fn panic(_info: &PanicInfo) -> ! {
    loop {}
}

#[no_mangle]
pub extern "C" fn main() -> ! {
    loop {
        // Main application loop
        cortex_m::asm::nop();
    }
}
EOF
\end{minted}

Add dependencies to \filepath{Cargo.toml}:\

\begin{minted}{toml}
[package]
name = "efr32mg24-minimal"
version = "0.1.0"
edition = "2021"

[profile.dev]
opt-level = "z"
debug = true
overflow-checks = false  # Critical for embedded

[profile.release]
opt-level = "z"
debug = true
overflow-checks = false
panic = "abort"
lto = true

codegen-units = 1

[dependencies]
cortex-m = "0.7"
cortex-m-rt = "0.7"

# Development dependencies
[dev-dependencies]
cortex-m-semihosting = "0.5"
\end{minted}

\subsection{Reference: Complete Project Example}
\label{subsec:reference-complete-project}

For a comprehensive step-by-step guide with all project files (build.rs, Cargo.toml, memory.x, etc.), see Section \ref{sec:complete-project-example}.

This section focuses on understanding the build process and analyzing the resulting binary.

\subsection{Understanding the Build Process}
\label{sec:understanding-build}

When you run \codeinline{cargo build --release}, here's what happens:
\begin{remark}
    This diagram shows the build pipeline for embedded Rust:
\end{remark}

\begin{figure}[H]
    \centering
    \begin{tikzpicture}[
        node/.style={rectangle, draw=black, thick, minimum width=2.5cm, minimum height=1cm},
        arrow/.style={-Latex, thick}
    ]

        % Nodes
        \node[node, fill=rustblue!20] (source) at (0, 0) {Source Code};
        \node[node, fill=yellow!20] (rustc) at (3, 0) {rustc};
        \node[node, fill=orange!20] (llvm) at (6, 0) {LLVM};
        \node[node, fill=green!20] (obj) at (9, 0) {Object File};
        \node[node, fill=red!20] (linker) at (9, -2) {Linker};
        \node[node, fill=purple!20] (bin) at (6, -2) {Binary};

        % External files
        \node[font=\small] (linker-script) at (9, 0.8) {linker.ld};
        \node[font=\small] (startup) at (9, 1.2) {startup_s};

        % Arrows
        \draw[arrow] (source) -- (rustc);
        \draw[arrow] (rustc) -- (llvm);
        \draw[arrow] (llvm) -- (obj);
        \draw[arrow] (obj) -- (linker);
        \draw[arrow] (linker) -- (bin);
        \draw[arrow, dashed] (linker-script) -- (linker);
        \draw[arrow, dashed] (startup) -- (linker);

        % Labels
        \node[font=\small, below=0.2cm of source] {Human};
        \node[font=\small, below=0.2cm of rustc] {Targeting};
        \node[font=\small, below=0.2cm of llvm] {ARM};
        \node[font=\small, below=0.2cm of obj] {Cortex-M33};
        \node[font=\small, below=0.2cm of bin] {thumbv8m};
    \end{tikzpicture}
    \caption{Rust embedded build pipeline}
    \label{fig:build-pipeline}
\end{figure}

\subsection{Compile Step}
\begin{minted}{bash}
$ cargo build --release --example minimal
   Compiling autocfg v1.1.0
   Compiling proc-macro2 v1.0.56
   Compiling cortex-m v0.7.7
   Compiling quote v1.0.27
   Compiling cortex-m-rt v0.7.3
   Compiling efr32mg24-minimal v0.1.0
    Finished release [optimized] target(s) in 1.24s
\end{minted}

Notice that cargo automatically downloads and compiles dependencies. For embedded,
this includes startup code and linker scripts via \codeinline{cortex-m-rt}.

\subsection{Binary Analysis}
\label{sec:binary-analysis}

Once built, let's examine the binary:

\subsection{Size Analysis}
\begin{minted}{bash}
$ cargo size --release --example minimal
   text    data     bss     dec     hex filename
     80       0       0      80      50 minimal
\end{minted}

What do these sections mean?
\begin{itemize}
    \item \textbf{text}: The actual machine code (read-only)
    \item \textbf{data}: Initialized static variables
    \item \textbf{bss}: Uninitialized static variables (zeroed at startup)
    \item \textbf{dec}: Total bytes (decimal)
    \item \textbf{hex}: Total bytes (hexadecimal)
\end{itemize}

Our minimal program is 80 bytes total. Let's see what's inside:

\subsection{Disassembly}
\begin{minted}{bash}
$ cargo objdump --release --example minimal -- -d
\end{minted}

Output (simplified):
\begin{minted}{text}
00000000 <main>:
       0: b580            push    {r7, lr}
       2: f000 f802       bl      0xa <_ZN4core4...
       6: bf00            nop
       8: e7fe            b.n     0x8 <main+0x8>

0000000a <HardFault>:
       a: e7fe            b.n     0xa <HardFault>
\end{minted}

Interpretation:
\begin{itemize}
    \item \codeinline{push {r7, lr}}: Save registers
    \item \codeinline{bl}: Branch with link (call function)
    \item \codeinline{nop}: No operation (our loop body)
    \item \codeinline{b.n}: Branch to self (infinite loop)
    \item \codeinline{HardFault}: Default exception handler
\end{itemize}

\subsection{Listing Sections}
\begin{minted}{bash}
$ cargo objdump --release --example minimal -- -h
\end{minted}

Shows memory layout:
\begin{minted}{text}
Idx Name          Size      VMA       LMA       File off  Algn
  0 .text         00000014  00000000  00000000  00000034  2**2
                  CONTENTS, ALLOC, LOAD, READONLY, CODE
  1 .data         00000000  20000000  20000000  00000034  2**0
                  CONTENTS, ALLOC, LOAD, DATA
  2 .bss          00000000  20000000  20000000  00000034  2**0
                  ALLOC
\end{minted}

\textbf{Explanation}:
\begin{itemize}
    \item \codeinline{.text} starts at 0x00000000 (flash memory)
    \item \codeinline{.data/.bss} start at 0x20000000 (SRAM)
    \item Flash: 1,536 KB on EFR32MG24
    \item SRAM: 256 KB on EFR32MG24
\end{itemize}

\subsection{Build Optimization Comparison}
\label{sec:optimization-comparsion}

Build with different optimization levels:

\begin{table}[H]
    \centering
    \begin{tabular}{l c c c c}
        \toprule
        \textbf{Profile} & \textbf{text} & \textbf{data} & \textbf{bss} & \textbf{Total} \\
        \midrule
        dev (opt-level='z') & 124 & 0 & 0 & 124 \\
        release (opt-level='z') & 80 & 0 & 0 & 80 \\
        release (opt-level='s') & 76 & 0 & 0 & 76 \\
        release (opt-level=3) & 68 & 0 & 0 & 68 \\
        \bottomrule
    \end{tabular}
    \caption{Size comparison by optimization level}
    \label{tab:optimization-size}
\end{table}

Optimization levels:
\begin{itemize}
    \item \codeinline{'z'}: Optimize for size aggressively
    \item \codeinline{'s'}: Optimize for size
    \item \codeinline{'3'}: Optimize for speed (may reduce size too)
\end{itemize}

Even the unoptimized debug build is only 124 bytes!

\subsection{Linking and Linker Scripts}
\label{sec:linking}

The linker script is crucial in embedded systems. It tells the linker where to place code and data in memory.
\subsection{Memory Layout}
Create \filepath{memory.x}:
\begin{minted}{text}
MEMORY
{
  /* Flash memory for code and constants */
  FLASH (rx) : ORIGIN = 0x00000000, LENGTH = 1536K

  /* RAM for variables and stack */
  RAM (rwx)  : ORIGIN = 0x20000000, LENGTH = 256K
}
\end{minted}

This matches the EFR32MG24B220 memory map:
\begin{itemize}
    \item Flash: 1,536 KB starting at 0x0000_0000
    \item SRAM: 256 KB starting at 0x2000_0000
\end{itemize}
\subsection{Examining the Linker Output}
\begin{minted}{bash}
$ cargo build --release --example minimal --verbose 2>&1 | grep "LINK"
\end{minted}

Shows the linker command:
\begin{minted}{text}
"rust-lld" "-flavor" "gnu"
  "-Tlink.x"
  "-Tmemory.x"
  "-o" ".../target/release/examples/minimal"
  ...
\end{minted}

\textbf{Key flags}:
\begin{itemize}
    \item \codeinline{rust-lld}: LLVM linker for embedded targets
    \item \codeinline{-Tlink.x}: Standard Cortex-M linker script
    \item \codeinline{-Tmemory.x}: Our custom memory layout
\end{itemize}

\subsection{Creating a Flashable Binary}
\label{sec:flashable-binary}

The ELF binary contains debug symbols and metadata. For flashing, we need a raw binary:
\subsection{Convert to Binary Format}
\begin{minted}{bash}
$ cargo objcopy --release --example minimal -- -O binary minimal.bin
\end{minted}

Check sizes:
\begin{minted}{bash}
$ ls -lh target/thumbv8m.main-none-eabihf/release/
-rwxr-xr-x  80K minimal          # ELF with symbols
-rwxr-xr-x  80B minimal.bin     # Raw binary for flash
\end{minted}

The binary is much smaller because it contains only the raw machine code without debug symbols, section headers, etc.
\subsection{Hex Format for Some Tools}
Some tools prefer Intel hex format:
\begin{minted}{bash}
$ cargo objcopy --release --example minimal -- -O ihex minimal.hex
\end{minted}

\subsection{Automated Build Script}
\label{sec:automated-build}

Create a rust file, \filepath{build.rs}, in your project root:
\begin{minted}{rust}
use std::env;
use std::fs;
use std::path::PathBuf;

fn main() {
    // Tell cargo to rerun build script if memory.x changes
    println!("cargo:rerun-if-changed=memory.x");

    // Copy memory.x to output directory
    let out = &PathBuf::from(env::var_os("OUT_DIR").unwrap());
    fs::copy("memory.x", out.join("memory.x")).unwrap();

    println!("cargo:rustc-link-search={}", out.display());
}
\end{minted}

This ensures the linker script is available during linking.

\subsection{Build Automation with Justfile}
\label{sec:justfile}

For complex projects, use \structure{just} (a command runner):
\begin{minted}{text}
# Justfile - Build automation

# Build release binary
build:
    cargo build --release --example minimal

# Show binary size
size:
    cargo size --release --example minimal

# Disassemble
disasm:
    cargo objdump --release --example minimal -- -d

# Show section info
sections:
    cargo objdump --release --example minimal -- -h

# Convert to binary
binary:
    cargo objcopy --release --example minimal -- -O binary minimal.bin

# Run all analysis
analyze: build size disasm sections binary
    @echo "Analysis complete!"
\end{minted}

Usage:
\begin{minted}{bash}
$ just build
$ just size
$ just analyze
\end{minted}
\subsection{Key Takeaways}
\label{sec:ch1-key-takeaways}

\begin{keytakeaways}
    \item Embedded Rust binaries are incredibly small (our minimal example is 80 bytes)
    \item Optimization settings significantly impact size (from 124 to 68 bytes)
    \item The build process involves: compile → link → optionally convert format
    \item Linker scripts control memory layout and are critical for embedded
    \item Cargo provides powerful tools for analyzing binaries (size, objdump, objcopy)
    \item The ELF format contains metadata; for flashing, use binary or hex format
    \item Scripting your build workflow improves productivity
\end{keytakeaways}

\begin{exercise}[Easy]
    Create a project template with a complete build script that:
    \begin{itemize}
        \item Builds release binary
        \item Shows size information
        \item Creates both .bin and .hex files
        \item Runs objdump analysis
    \end{itemize}
\end{exercise}

\begin{exercise}[Medium]
    Compare binaries built with different optimization levels:
    \begin{itemize}
        \item What differences do you see in objdump output?
        \item How much code can the linker remove from our minimal example?
        \item Where does libc start inserting startup code?
    \end{itemize}
\end{exercise}

\begin{exercise}[Hard]
    Write a custom linker script that:
    \begin{itemize}
        \item Places the vector table at a specific address
        \item Creates a separate section for time-critical code
        \item Reserves space at the end of RAM for a stack canary
        \item Adds checksums for firmware verification
    \end{itemize}
    Then verify it works with objdump.
\end{exercise}


% ============================================================================
% SECTION 01.07: What Goes In Each Crate, and Why
% ============================================================================
\section{What Goes In Each Crate, and Why}
\label{sec:what-goes-in-each-crate}

The PAC/HAL/BSP pattern from Section~\ref{sec:pac-hal-bsp-architecture} says the layers exist.
It does not say where one ends and the next begins, and every project fills that gap
differently. That gap is what you face the first time you add support for a new MCU or a new
board — which is the situation this book is written for, and the one it should leave you
equipped to handle.

This section argues the choice, compares it with what four mature Rust embedded ecosystems do,
gives a decision procedure you can apply to hardware that is not in this book, and states the
rules that keep the answer from drifting.

\subsection{The Question, Stated Sharply}

Three pieces of code make the ambiguity visible. All three are real, from this project:

\begin{enumerate}
    \item \textbf{A pin pair.} The IMU on this board answers on \texttt{SDA = PB2},
          \texttt{SCL = PB3} — the opposite of what the vendor's pin map declares. Whose
          knowledge is that?
    \item \textbf{A power gate.} The same sensor does not answer until \texttt{PD5} is driven
          high. Is that a driver detail or a board detail?
    \item \textbf{An algorithm.} A fall classifier takes acceleration samples and decides
          whether the device fell. It has no registers, no pins and no clock. Where does it
          live so that it stays testable?
\end{enumerate}

A HAL that answers (1) and (2) knows too much: every new board becomes a fork of the HAL. An
example that answers them re-derives them, and the same investigation gets repeated by the
next person — the failure this project has already paid for once.

\subsection{Three Candidate Layouts}

\paragraph{Layout A — no BSP at all; wiring lives in the examples.}
This is the smallest possible structure: the HAL offers drivers, and each example says which
pins to use. It is what the STM32 and ESP32 ecosystems do, and it is defensible when the point
of the example is to teach the peripheral.

Its cost appears the moment the wiring is not obvious. This board's IMU needed a swapped
\texttt{SDA}/\texttt{SCL} and a power gate before it answered anything; putting that in every
example means every example can get it wrong independently, and the knowledge is only as good
as the most recently written one.

\paragraph{Layout B — one BSP crate, features for the optional hardware.}
A single \codeinline{xiao-}\allowbreak\codeinline{mg24-bsp} with features
\codeinline{sense} and \codeinline{expansion-base}. It looks like the cheapest correct answer, and it has a specific
trap that the nRF ecosystem documents in its own repository: \emph{cargo unifies features
across a workspace}, so a build that enables the base for one target can silently enable it
for another. More importantly, the API lies: a \codeinline{ExpansionBase} type that exists in a
build without the base has to answer ``no OLED'' at runtime, and every method that could have
been a compile error becomes an \codeinline{Option}.

\paragraph{Layout C — one BSP crate per board, and a feature for a variant of one board.}
The module gets its own crate (LED on both variants; IMU and microphone behind
\codeinline{sense}), the Expansion Base gets its own crate that depends on the module's, and
the dependency arrow points one way. A user who does not own the base never compiles it, and
never sees an \codeinline{Option} that only exists because two boards were welded together.

\subsection{What Four Ecosystems Actually Do}

Read from the repositories, not from memory:

\begin{table}[H]
\centering
\footnotesize
\begin{tabular}{@{}>{\raggedright\arraybackslash}p{2.3cm}>{\raggedright\arraybackslash}p{2.6cm}>{\raggedright\arraybackslash}p{2.4cm}>{\raggedright\arraybackslash}p{3.4cm}>{\raggedright\arraybackslash}p{3.0cm}@{}}
\hline
\textbf{ecosystem} & \textbf{chip selection} & \textbf{BSP crates} & \textbf{where examples live} & \textbf{where wiring is written} \\
\hline
\textbf{rp-hal} (RP2040/RP2350) & one HAL per chip & \textbf{yes}, one per board,
  in a \emph{separate repository} (\codeinline{rp-hal-boards/boards/}) & in each BSP, plus
  separate example crates & in the BSP, once \\
\textbf{nrf-hal} & one HAL per chip (\codeinline{nrf52840-hal}, \codeinline{nrf9160-hal}, \dots)
  + \codeinline{nrf-hal-common} & no & in the HAL repo, board chosen by feature & in the
  examples \\
\textbf{stm32f4xx-hal} & one HAL per family & no & in the HAL crate &
  in the examples (pin numbers inline) \\
\textbf{esp-hal} & one crate, chip by feature & no & one \codeinline{examples/} grouped by
  capability & in the examples \\
\hline
\end{tabular}
\caption{How four Rust embedded ecosystems split PAC, HAL, BSP and examples. Verified from the
repositories in November 2026.}
\label{tab:ecosystem-layouts}
\end{table}

Two details in that table are worth more than the table.

First, \textbf{rp-hal puts its BSPs in a different repository from the HAL}. That is stronger
than a folder: the HAL's continuous integration cannot depend on a board, because the board is
not present. The BSP crates describe themselves as a HAL for one board that re-exports the
chip HAL — a facade, exactly as this book calls it.

Second, \textbf{nrf-hal ships a workspace template instead of a workspace}, and says why:
\emph{``Since Cargo resolves all features for all crates in the workspace together, it is not
possible to have a workspace with all the crates together, since they enable incompatible
features.''} That sentence is the argument against Layout B in the ecosystem's own words.

\subsection{The Decision, and the Rule Behind Each Part}

\begin{table}[H]
\centering
\footnotesize
\begin{tabular}{@{}>{\raggedright\arraybackslash}p{3.4cm}>{\raggedright\arraybackslash}p{5.2cm}>{\raggedright\arraybackslash}p{5.2cm}@{}}
\hline
\textbf{crate / folder} & \textbf{what goes in} & \textbf{what must never go in} \\
\hline
\codeinline{efr32mg24-pac} & registers, generated from the SVD & anything hand-written that
  a generator could produce \\
\codeinline{efr32mg24-hal} & drivers: which bytes, in which order, with which timing & pins,
  addresses, power gates, board names \\
\codeinline{xiao-mg24-bsp} & the module: LED; behind \codeinline{sense}, the IMU bus, its
  power gate and the microphone & any device that is not on the module \\
\codeinline{xiao-expansion-base-bsp} & the base: OLED, BME68x, RTC, buzzer, button, Grove,
  microSD, and the base's own bus & anything that exists without the base \\
\codeinline{libraries/} & domain policy with no hardware at all & any dependency on the
  hardware stack, direct or transitive \\
\hline
\end{tabular}
\caption{Ownership. The right-hand column is what the gate checks; the left two are what a
review checks.}
\label{tab:crate-ownership}
\end{table}

Three rules follow, and each one exists because of a specific failure:

\begin{keytakeaways}
    \item \textbf{A BSP describes one piece of plastic.} Two boards that stack are two
          crates, because a user may own only one of them.
    \item \textbf{A variant of one board is a feature, not a crate.} The Sense module is the
          same PCB with more parts fitted; a separate crate would duplicate the board init
          and the LED to add two sensors.
    \item \textbf{A dependency arrow never points at a board.} The base depends on the
          module; the module does not know the base exists. This is the property that makes
          ``does it work without the base?'' answerable by the compiler instead of by a test
          matrix.
\end{keytakeaways}

\subsection{Where the Algorithm Goes}

The classifier deserves its own sentence because it is where layering usually breaks. It has
no hardware, but the temptation is to put it ``near'' the driver that produces its input — in
the same crate as the IMU, or in the example that uses it.

Both are wrong for the same reason: the moment it shares a crate with anything that needs a
target triple, its tests need a target triple, and the tests are what keep it honest. In this
project it is a crate with \textbf{zero dependencies}, and the gate fails if its dependency
tree ever contains a hardware crate:

\begin{minted}{bash}
$ cargo tree -p motion-detection | grep efr32mg24
$ echo $?
1        # nothing from the hardware stack: the rule holds
\end{minted}

The contract is worth quoting, because it is the whole testability argument in one line:

\begin{minted}{rust}
// [i32; 3] of milligravity, and the interval between samples in microseconds.
pub fn update_us(&mut self, accel_mg: [i32; 3], dt_us: u32) -> Option<Event>
\end{minted}

No sensor type, no bus, no clock. Which is why nine tests run in under a second on a laptop,
and why the same code can be replayed against a recording at ten times real speed.

\subsection{And Where the Examples Go}

The layout of the examples follows the same principle as the crates, with one addition for
hardware that stacks: \textbf{an example lives in the crate of the largest hardware
requirement it has, and states that requirement in its first comment.}

\begin{minted}{text}
crates/xiao-mg24-bsp/examples/            runs with the module alone on the desk
crates/xiao-expansion-base-bsp/examples/  needs the module (Sense) and the base
crates/efr32mg24-pac/examples/            chip level: any pin is legitimate
\end{minted}

The alternative — every example in one folder — looks tidier and is worse: the crate of the
module would have no examples, and an API that nothing builds is an API that nothing checks.
A crate with hardware features and zero examples should fail a gate, not pass it quietly.

\subsection{A Decision Procedure for Hardware Not In This Book}

The four questions below decide the layout. They take minutes and they are the difference
between a structure that survives the second board and one that has to be refactored.

\begin{enumerate}
    \item \textbf{How many chips?} If more than one, the HAL is per chip (or per family, when
          the parts are close enough that one driver set serves them). This is what
          \codeinline{nrf-hal} does per chip and \codeinline{stm32f4xx-hal} does per family;
          \codeinline{esp-hal} instead keeps one crate and selects the chip by feature, which
          works when the vendor's crates already share their register code.
    \item \textbf{How many boards, and can one sit on top of another?} Boards that are
          alternatives become one BSP crate each. Boards that \emph{stack} become one crate
          each \emph{plus} a one-way dependency: the carrier depends on the module. A single
          crate that knows both cannot answer ``does this work without the carrier?''.
    \item \textbf{Are there variants of the same board?} A variant — same PCB, different
          parts fitted — is a \emph{feature}, not a crate. A separate crate duplicates the
          board init and everything the variants share, in exchange for a name that repeats
          what the feature flag already says.
    \item \textbf{Is the wiring obvious or hard-won?} Obvious wiring (an LED on a documented
          pin) may live in an example. Hard-won wiring — a swapped bus, a power gate, a pin
          that the vendor's own map gets wrong — belongs in a BSP, once, with the measurement
          recorded next to it. This is the question that decides between Layout A and
          Layout C for most hobby-scale projects.
\end{enumerate}

The same questions decide where an algorithm goes, with one addition: \textbf{if it can be
tested on a host, it must be}, so it cannot share a crate with anything that needs a target
triple.

\subsection{Adding Support for a New Chip or Board: the Order of Work}

Having decided the layout, the order below is what this project would follow again. Each step
produces something checkable before the next one starts.

\begin{enumerate}
    \item \textbf{Registers first.} Obtain the SVD, generate the PAC, and build one example
          that only writes a register and reads it back. If the SVD is wrong, everything
          downstream inherits the error, and a five-line example finds that in minutes.
    \item \textbf{Clock and pins.} Prove that the peripheral clock gate is on and that a GPIO
          changes state. This is the step where a missing
          \codeinline{CMU\_CLKEN0} bit appears as a bus fault at the first register access.
    \item \textbf{One transport.} Bring up the simplest bus the board has, and read a device
          identifier through it. A chip identifier or a WHO\_AM\_I is a free proof that
          addressing, timing and the byte order are all correct.
    \item \textbf{Drivers, in the HAL, with pins as arguments.} Nothing in this layer may name
          a port. If a driver needs a pin, take it as a parameter; the BSP will supply it.
    \item \textbf{The BSP for one board, and its examples.} Put the wiring in one place, in
          the crate named after the board whose documentation describes the feature, and
          write the example that proves the board works end to end.
    \item \textbf{A second board, if there is one, and the dependency direction.} Add the
          carrier crate only after the module crate stands alone, and make the dependency
          point one way. If it cannot, the split is wrong.
    \item \textbf{The gate.} Before more examples are written, make the machine check what
          the review would check: zero warnings, every example builds, every domain crate
          tests on the host, and no domain crate depends on the hardware stack.
\end{enumerate}

\begin{warningbox}
    \textbf{The step people skip is the first one.} Generating the PAC and immediately
    writing drivers feels productive, and it hides an SVD error until the driver is
    complicated enough that the error looks like a driver bug. The five-line register
    read-back is the cheapest test in the whole project.
\end{warningbox}

\subsection{What This Costs}

Honesty about a design decision includes its price:

\begin{itemize}
    \item \textbf{More manifests.} Two BSP crates mean two \filepath{Cargo.toml} files, two
          example folders and a dependency between them that has to be kept one-way.
    \item \textbf{Some knowledge is duplicated on purpose.} The clock setup lives in the
          module crate and the base depends on the module rather than re-initialising; a
          project with two unrelated boards could not share it this way.
    \item \textbf{The split is a judgement, not a theorem.} ``Which plastic is this pin on?''
          is unambiguous for an LED and less so for a Grove connector that reaches the base
          through the module's castellated pads. The rule that resolves it: the BSP is named
          after the board whose \emph{documentation} describes the feature.
\end{itemize}

\begin{notebox}
    \textbf{The test to apply to your own board.} Ask which of your examples would still make
    sense to someone holding only the chip. If the answer is ``none of them'', the board has
    swallowed the crate and a split is due. If the answer is ``all of them'', you may not need
    a BSP at all — Layout A is not a mistake, it is a choice with a cost that is sometimes
    worth paying.
\end{notebox}


\clearpage

% ============================================================================
% Chapter Summary and Key Takeaways
% ============================================================================

\begin{keytakeaways}
    \item Rust provides memory safety without sacrificing performance through its ownership system
    \item Embedded Rust programs use \codeinline{no_std} and \codeinline{no_main} for bare-metal operation
    \item The PAC/HAL/BSP architecture provides clean separation of concerns and promotes code reuse
    \item Cross-compilation requires the correct target specification (thumbv8m.main-none-eabihf for EFR32MG24)
    \item Building embedded binaries involves multiple steps: compile, link, and optionally convert to binary/hex format
    \item All necessary tools (rustup, probe-rs, cargo-binutils) are available cross-platform
    \item Our minimal program compiles to just 80 bytes in release mode
\end{keytakeaways}

